Los experimentos fueron ejecutados en un equipo con procesador \textbf{Intel\textsuperscript{\textregistered} Core\texttrademark{} i5-11400H de 11\textsuperscript{a} generación a 2,70 GHz (2,69 GHz efectivos)}, con \textbf{16,0 GB de RAM} (15,8 GB utilizables). El sistema operativo utilizado fue \textbf{Windows 11 / Ubuntu 24.04.4 LTS}, y el compilador empleado fue \textbf{gcc (Ubuntu 13.3.0-6ubuntu2\textasciitilde 24.04.1) versión 13.3.0}.

Se utilizaron los conjuntos de datos descritos en el enunciado de la tarea. Para el problema de ordenamiento, se emplearon arreglos de tamaño $n \in \{10^1, 10^3, 10^5, 10^7\}$ con variaciones en su tipo (ascendente, descendente, aleatorio) y dominio ($D1, D7$). Para la multiplicación de matrices, se usaron matrices cuadradas de dimensiones $n \in \{2^4, 2^8, 2^{10}\}$ con estructuras dispersas, diagonales y densas, bajo los dominios $D0$ y $D10$. En ambos casos, se consideraron tres muestras aleatorias ($a, b, c$) por configuración para las mediciones experimentales.